\section*{Chapitre \ref{ch:g8:speed-of-light} --- La vitesse de la lumière~; les distances dans l'Univers}

\begin{solution}{exo:g8:speed-of-light:1}
$c = \qty{300000}{km/s} = 3 \times 10^{8}\ \unit{m/s}$. Les lanternes
des collines n'ont mesuré que les réflexes humains~: la \omterm{def:g1:light-and-shadows:source}{lumière}
franchissait les collines en millionièmes de seconde, bien en deçà de
tout chronométrage à la main.
\end{solution}

\begin{solution}{exo:g8:speed-of-light:2}
Le satellite de Jupiter s'éclipsait selon un horaire ponctuel --- une
horloge du ciel. La \emph{nouvelle} de chaque éclipse arrivait en
retard chaque fois que la Terre se tenait loin de Jupiter, en avance
quand elle en était proche. Le retard mesurait le temps de trajet de
la \omterm{def:g1:light-and-shadows:source}{lumière} sur la largeur de l'\omterm{def:g6:sun-earth-moon:axisorbit}{orbite} terrestre --- et donc la \omterm{def:g7:motion-average-speed:formula}{vitesse}
de la \omterm{def:g1:light-and-shadows:source}{lumière}.
\end{solution}

\begin{solution}{exo:g8:speed-of-light:3}
\omterm{def:g2:moon-in-the-sky:moon}{Lune}~: $384000 \div 300000 \approx \qty{1.3}{s}$. Soleil~:
$150000000 \div 300000 = \qty{500}{s}$ --- huit minutes vingt. La
\omterm{def:g1:light-and-shadows:source}{lumière} du Soleil montre toujours le Soleil tel qu'il était il y a
huit minutes~: une vieille nouvelle par construction.
\end{solution}

\begin{solution}{exo:g8:speed-of-light:4}
Une distance~: le chemin que la \omterm{def:g1:light-and-shadows:source}{lumière} parcourt en un an --- environ
$9.5 \times 10^{12}$ km.
\end{solution}

\begin{solution}{exo:g8:speed-of-light:5}
Quatre ans. Les réponses varient --- il y a quatre ans, tu en étais à
plusieurs niveaux en arrière dans ce livre même.
\end{solution}

\begin{solution}{exo:g8:speed-of-light:6}
À onze minutes environ à l'aller, une correction de \omterm{def:g6:describing-motion:trajectory}{trajectoire} répond
à une \omterm{prop:g5:mirrors-reflection:image}{image} vieille d'au moins vingt-deux minutes --- le rover serait
dans la crevasse avant l'arrivée du moindre frémissement de manche.
D'où des plans, et non des manches à balai.
\end{solution}

\begin{solution}{exo:g8:speed-of-light:7}
\omterm{def:g2:moon-in-the-sky:moon}{Lune}~: environ \qty{1.3}{s}. Soleil~: huit minutes. \omterm{def:g4:solar-system:star}{Étoile} la plus
proche~: quatre ans. Galaxie voisine~: deux millions et demi
d'années.
\end{solution}

\begin{solution}{exo:g8:speed-of-light:8}
La \omterm{def:g1:light-and-shadows:source}{lumière}. La \omterm{def:g7:motion-average-speed:formula}{vitesse} $c$ s'est révélée plus stable que toute barre
de métal ou règle nationale --- si bien que la règle se déduit
désormais de la \omterm{def:g1:light-and-shadows:source}{lumière}, et non la \omterm{def:g1:light-and-shadows:source}{lumière} d'une mesure à la règle.
\end{solution}

\begin{solution}{exo:g8:speed-of-light:9}
Aller-retour~: $300000 \times 2.6 = 780000$ km~; la \omterm{def:g2:moon-in-the-sky:moon}{Lune} se tient à
la moitié --- \qty{390000}{km}. La règle du sonar se transpose tout
entière~: tout \omterm{ex:g4:how-sound-travels:echo}{écho} paie l'aller-retour, diviser par deux avant de
croire.
\end{solution}

\begin{solution}{exo:g8:speed-of-light:10}
Lampe~: $3 \div (3 \times 10^{8}) = 10^{-8}$ s --- un centième de
millionième de seconde~: même la chambre est une archive, trop peu
profonde pour se remarquer. Puis le Soleil (huit minutes), l'\omterm{def:g4:solar-system:star}{étoile}
(un siècle), la galaxie (deux millions et demi d'années) --- plus tu
regardes loin, plus la page est ancienne.
\end{solution}

\begin{solution}{exo:g8:speed-of-light:11}
«~Mardi dernier, la \omterm{def:g1:light-and-shadows:source}{lumière} de l'explosion d'une \omterm{def:g4:solar-system:star}{étoile} \emph{nous est
parvenue} --- l'explosion elle-même a eu lieu il y a quelque $8000$
ans, et sa nouvelle voyage depuis lors.~»
\end{solution}

\begin{solution}{exo:g8:speed-of-light:12}
$3 \times 10^{8}\ \text{km} \div 1000\ \text{s} = 3 \times 10^{5}$
\unit{km/s} --- la valeur moderne pile (l'estimation historique, avec
des chiffres d'\omterm{def:g6:sun-earth-moon:axisorbit}{orbite} plus grossiers, tombait à environ un quart
près~: pour les années seize cents, une première réponse
stupéfiante).
\end{solution}

\begin{solution}{pb:g8:speed-of-light:1}
\textbf{1.} Environ \qty{1.3}{s} après que tu as parlé~; l'accusé de
réception le plus rapide revient après l'aller-retour, soit environ
\qty{2.6}{s}.
\textbf{2.} $36000 \div 300000 = \qty{0.12}{s}$ à l'aller~; une
question et sa réponse montent et descendent deux fois --- près d'une
demi-seconde d'hésitation, juste au seuil de l'oreille.
\textbf{3.} «~L'éruption a jailli huit minutes vingt avant que nos
écrans ne s'allument --- nous la regardons telle qu'elle
\emph{était}.~»
\textbf{4.} La liaison lunaire~: \qty{1.3}{s} contre \qty{0.12}{s}
--- environ dix fois la marge.
\textbf{5.} $1.8 \times 10^{8} \div (3 \times 10^{5}) = 600$ s~: dix
minutes à l'aller.
\textbf{6.} L'aller-retour dure vingt minutes~: l'\omterm{prop:g5:mirrors-reflection:image}{image} du sauvetage
a dix minutes et l'ordre de sauvetage arrive dix minutes trop tard ---
la crevasse gagne avec dix-neuf minutes d'avance. Le rover doit donc
emporter ses propres réflexes~: un arrêt d'urgence embarqué.
\textbf{7.} Le dernier morceau doit arriver pour 21:00~; il quitte la
Terre à $21{:}00 - 10$ min $= 20{:}50$~; la transmission a commencé
trois minutes plus tôt~: départ à 20:47.
\textbf{8.} Le retard double, à environ vingt et une minutes à
l'aller~: les \omterm{def:g4:solar-system:planet}{planètes} n'ont pas «~une~» distance --- les deux mondes
sont en \omterm{def:g6:sun-earth-moon:axisorbit}{orbite}, et toute conversation se planifie contre une cible
mouvante.
\textbf{9.} $300000 \times 86400 \approx 2.6 \times 10^{10}$ km ---
vingt-six mille millions de kilomètres~; une question et sa réponse
prennent deux jours pleins.
\textbf{10.} La \omterm{def:g1:light-and-shadows:source}{lumière} atteint l'\omterm{def:g4:solar-system:star}{étoile} en environ quatre ans~; la
réponse la plus rapide se pose environ huit ans après l'émission de ce
soir.
\textbf{11.} «~Cette \omterm{def:g1:light-and-shadows:source}{lumière} a quitté sa galaxie il y a deux millions
et demi d'années~: l'instantané rediffuse la Terre des tout premiers
outils de pierre --- notre page la plus profonde de l'archive à
l'œil nu.~»
\textbf{12.} Par exemple~: «~Rien n'est instantané --- tout signal, le
nôtre compris, voyage au mieux à $c$. Les manches à balai perdent
contre le retard au-delà de la \omterm{def:g2:moon-in-the-sky:moon}{Lune}~: les machines lointaines doivent
penser par elles-mêmes. Et tout télescope est un lecteur d'archives
--- plus l'objet est loin, plus la nouvelle est ancienne, et aucune
édition plus fraîche n'est imprimable.~»
\end{solution}
